% ═══════════════════════════════════════════════════════════════
% FORSCHERPROTOKOLL: Elektromotor-Experiment (Kl. 9)
% A4, mit AFB III – SYNCHRONISIERT MIT KARTE (40 Min)
% ═══════════════════════════════════════════════════════════════

\documentclass[10pt, a4paper]{article}

\usepackage[utf8]{inputenc}
\usepackage[T1]{fontenc}
\usepackage[ngerman]{babel}

% --- SCHRIFTEN ---
\usepackage{helvet}
\renewcommand{\familydefault}{\sfdefault}

\usepackage{microtype}
\usepackage[margin=1.3cm, top=1.3cm, bottom=1.3cm]{geometry}
\usepackage{enumitem}
\usepackage{tabularx}
\usepackage{booktabs}
\usepackage{array}
\usepackage{multicol}
\usepackage{tikz}
\usepackage{amssymb}

\usepackage{fancyhdr}
\usepackage{lastpage}

\usepackage{xcolor}
\usepackage{tcolorbox}
\tcbuselibrary{skins,breakable}

% ═══════════════════════════════════════════════════════════════
% FARBEN
% ═══════════════════════════════════════════════════════════════

\definecolor{physik}{HTML}{2c5aa0}
\definecolor{mittelgrau}{HTML}{888888}
\definecolor{warnung}{HTML}{E65100}
\definecolor{erfolg}{HTML}{2E7D32}
\definecolor{stern3}{HTML}{FFB300}

\colorlet{fachfarbe}{physik}

% ═══════════════════════════════════════════════════════════════
% VARIABLEN
% ═══════════════════════════════════════════════════════════════

\newcommand{\fach}{Physik}
\newcommand{\thema}{Elektromotor-Experiment}
\newcommand{\klasse}{9}

% ═══════════════════════════════════════════════════════════════
% KOPF- UND FUSSZEILE
% ═══════════════════════════════════════════════════════════════

\pagestyle{fancy}
\fancyhf{}
\renewcommand{\headrulewidth}{0.4pt}
\renewcommand{\footrulewidth}{0pt}
\lhead{\small \fach{} Klasse \klasse{} | \thema}
\rhead{\small Name: \underline{\hspace{3cm}}}
\cfoot{\small\textcolor{mittelgrau}{Seite \thepage{} von \pageref{LastPage}}}

% ═══════════════════════════════════════════════════════════════
% BOXEN – NUR RAHMEN, KEINE FARBFÜLLUNG
% ═══════════════════════════════════════════════════════════════

\newtcolorbox{aufgabenbox}[2][]{
    colback=white,
    colframe=fachfarbe,
    coltitle=black,
    colbacktitle=white,
    boxrule=1pt,
    arc=2pt,
    left=5pt, right=5pt, top=3pt, bottom=3pt,
    before skip=5pt, after skip=4pt,
    fonttitle=\bfseries\small,
    title={#2},
    #1
}

\newtcolorbox{knobelbox}[1][]{
    colback=white,
    colframe=stern3!80!black,
    coltitle=black,
    colbacktitle=white,
    boxrule=1.5pt,
    arc=3pt,
    left=5pt, right=5pt, top=3pt, bottom=3pt,
    before skip=6pt, after skip=4pt,
    fonttitle=\bfseries\small,
    title=#1
}

\newtcolorbox{materialbox}{
    colback=white,
    colframe=fachfarbe,
    boxrule=1pt,
    arc=0pt,
    left=6pt, right=6pt, top=4pt, bottom=4pt,
    before skip=4pt, after skip=4pt
}

\newtcolorbox{antwortbox}{
    colback=white,
    colframe=mittelgrau,
    boxrule=0.5pt,
    arc=0pt,
    left=3pt, right=3pt, top=2pt, bottom=2pt,
    before skip=2pt, after skip=2pt
}

\newtcolorbox{bonusbox}[1][]{
    colback=white,
    colframe=erfolg,
    coltitle=black,
    colbacktitle=white,
    boxrule=1pt,
    arc=2pt,
    left=4pt, right=4pt, top=3pt, bottom=3pt,
    fonttitle=\bfseries\small,
    title=#1
}

\newtcolorbox{extrabox}[1][]{
    colback=white,
    colframe=fachfarbe,
    coltitle=black,
    colbacktitle=white,
    boxrule=1pt,
    arc=2pt,
    left=4pt, right=4pt, top=3pt, bottom=3pt,
    fonttitle=\bfseries\small,
    title=#1
}

% ═══════════════════════════════════════════════════════════════
% BEFEHLE
% ═══════════════════════════════════════════════════════════════

\newcommand{\luecke}[1]{\underline{\hspace{#1}}}
\newcommand{\sternI}{$\star$}
\newcommand{\sternII}{$\star\star$}
\newcommand{\sternIII}{$\star\star\star$}

\setlength{\parindent}{0pt}
\setlength{\parskip}{1pt}

% ═══════════════════════════════════════════════════════════════
% DOKUMENT
% ═══════════════════════════════════════════════════════════════

\begin{document}

% ═══════════════════════════════════════════════════════════════
% HEADER
% ═══════════════════════════════════════════════════════════════

\begin{center}
    {\large\bfseries\textcolor{fachfarbe}{FORSCHERPROTOKOLL: Elektromotor}}
\end{center}
\vspace{-0.4em}
\hrule
\vspace{0.3em}

% ═══════════════════════════════════════════════════════════════
% MATERIAL
% ═══════════════════════════════════════════════════════════════

\begin{materialbox}
\small\textbf{Material:} Motor \textbullet{} Netzgerät \textbullet{} Kabel \textbullet{} Stabmagnet \textbullet{} Elektromagnet ~{\footnotesize(Aufbau/Abbau: siehe Experimentierkarte)}
\end{materialbox}

% ═══════════════════════════════════════════════════════════════
% ZWEI SPALTEN FÜR AUFGABEN 1-4
% ═══════════════════════════════════════════════════════════════

\begin{multicols}{2}

% AUFGABE 1
\begin{aufgabenbox}{Aufgabe 1: Motor starten \hfill \sternI}
\small
Bringt den Motor zum Drehen.

\textbf{Was braucht ihr mindestens?}

\begin{antwortbox}
\vspace{0.6cm}
\end{antwortbox}
\end{aufgabenbox}

% AUFGABE 2
\begin{aufgabenbox}{Aufgabe 2: Drehrichtung ändern \hfill \sternII}
\small
Ändert die Drehrichtung -- \textbf{ohne} den Motor umzudrehen.

\textbf{Zwei verschiedene Wege:}

\begin{antwortbox}
\textbf{1:} \luecke{4.5cm}

\vspace{0.35cm}

\textbf{2:} \luecke{4.5cm}
\end{antwortbox}
\end{aufgabenbox}

% AUFGABE 3
\begin{aufgabenbox}{Aufgabe 3: Magnet-Vergleich \hfill \sternII}
\small
Ersetzt den Stabmagneten durch den Elektromagneten.

\begin{antwortbox}
\begin{tabular}{@{}l@{\hspace{2mm}}l@{}}
\textbf{Stabmagnet:} & $\square$ ja $\square$ nein \\[0.25cm]
\textbf{Elektromagnet:} & $\square$ ja $\square$ nein \\
\end{tabular}

\vspace{0.15cm}
\textbf{Unterschied:}

\vspace{0.4cm}
\end{antwortbox}
\end{aufgabenbox}

\columnbreak

% AUFGABE 4
\begin{aufgabenbox}{Aufgabe 4: Spannung verändern \hfill \sternI}
\small
Dreht die Spannung langsam von 3\,V auf 12\,V.

\begin{antwortbox}
\textbf{Beobachtung:}

\vspace{0.6cm}
\end{antwortbox}
\end{aufgabenbox}

% AUFGABE 5
\begin{aufgabenbox}{Aufgabe 5: Magnet-Abstand \hfill \sternI}
\small
Verändert den Abstand zwischen Magnet und Motor.

\begin{antwortbox}
\textbf{Beobachtung:}

\vspace{0.6cm}
\end{antwortbox}
\end{aufgabenbox}

% AUFGABE 6
\begin{aufgabenbox}{Aufgabe 6: Motor stoppen \hfill \sternII}
\small
Bringt den Motor zum \textbf{Stillstand}, obwohl Strom fließt.

\begin{antwortbox}
\textbf{Wie?}

\vspace{0.6cm}
\end{antwortbox}
\end{aufgabenbox}

% AUFGABE 7
\begin{aufgabenbox}{Aufgabe 7: Lautstärke \hfill \sternII}
\small
Könnt ihr den Motor \textbf{lauter} oder \textbf{leiser} machen?

\begin{antwortbox}
\textbf{Was beeinflusst die Lautstärke?}

\vspace{0.5cm}
\end{antwortbox}
\end{aufgabenbox}

\end{multicols}

% ═══════════════════════════════════════════════════════════════
% KNOBELAUFGABEN (AFB III) - BEIDE PFLICHT
% ═══════════════════════════════════════════════════════════════

\begin{knobelbox}[KNOBELAUFGABEN \sternIII{} \hfill Bearbeitet BEIDE!]

\begin{multicols}{2}

\small
\textbf{A: Das Totpunkt-Problem}

Manchmal bleibt der Motor stehen, obwohl Strom fließt.

\textbf{Erkläre mit der Lorentzkraft}, warum das passiert. Zeichne eine Skizze.

\begin{antwortbox}
\textbf{Erklärung:}
\vspace{0.9cm}

\textbf{Skizze:}
\vspace{1.1cm}
\end{antwortbox}

\columnbreak

\small
\textbf{B: Magnet-Bewertung}

Für einen \textbf{Spielzeugmotor}: Stabmagnet oder Elektromagnet? Begründe!

\begin{antwortbox}
\begin{tabular}{@{}l@{ }l@{}}
\textbf{Stabmagnet} & \textbf{Elektromagnet} \\
+ \luecke{2.1cm} & + \luecke{2.1cm} \\[0.25cm]
$-$ \luecke{2.1cm} & $-$ \luecke{2.1cm} \\
\end{tabular}

\vspace{0.15cm}
\textbf{Besser:} \luecke{2.5cm}

\textbf{Weil:}
\vspace{0.4cm}
\end{antwortbox}

\end{multicols}

\end{knobelbox}

% ═══════════════════════════════════════════════════════════════
% BONUS + ÜBERRASCHUNG
% ═══════════════════════════════════════════════════════════════

\vspace{0.15cm}

\begin{multicols}{2}

\begin{bonusbox}[BONUS (für Schnelle)]
\small
Bei welcher Spannung startet der Motor \textbf{gerade so}?

\vspace{0.15cm}
\textbf{Startspannung:} \luecke{1.5cm} V
\end{bonusbox}

\columnbreak

\begin{extrabox}[FAZIT]
\small
Was hast du heute über den Motor herausgefunden?

\vspace{0.15cm}
\luecke{5.5cm}
\end{extrabox}

\end{multicols}

\vfill

\begin{center}
\scriptsize\textcolor{mittelgrau}{Forscherprotokoll | Physik Klasse 9 | 40 Min}
\end{center}

\end{document}
